\Article{38}
We now turn to the congruence of numbers with respect to complex moduli. At the outset of this investigation, however, we should indicate how the domain of complex quantities can be made visible to intuition.

Every real quantity can be expressed by a segment of a line infinite in both directions, taken from an arbitrary starting point and measured against an arbitrary segment adopted as the unit. It can thus be represented by the other endpoint, with points on one side of the starting point representing positive quantities and those on the other negative quantities. In the same way, any complex quantity can be represented by a point in an infinite plane in which a fixed line corresponds to the real quantities: the complex quantity $x+iy$ by the point with abscissa $=x$ and ordinate, taken positively on one side of the line of abscissas and negatively on the other, $=y$. In this way we may say that any complex quantity measures the difference between the position of the point to which it corresponds and that of the starting point. The positive unit denotes an arbitrary fixed displacement in an arbitrary fixed direction; the negative unit an equally large displacement in the opposite direction; and the imaginary units equally large displacements in the two perpendicular lateral directions.

This sheds considerable light on the metaphysics of the quantities we call imaginary. If the starting point is denoted by $(0)$, and two complex quantities $m,m'$ correspond to points $M,M'$, whose positions relative to $(0)$ they express, the difference $m-m'$ is nothing other than the position of $M$ relative to $M'$. On the other hand, if the product $mm'$ represents the position of $N$ relative to $(0)$, it is easy to see that this position is determined by the position of $M$ relative to $(0)$ in the same way as the position of $M'$ is determined by the position of the point corresponding to the positive unit. Thus we may aptly say that the positions of the points corresponding to the complex quantities $mm'$
\SourceLeaf{120}%
$m,m',1$ form a \emph{proportion}. But we reserve a fuller treatment of this matter for another occasion. The difficulties in which the theory of imaginary quantities is thought to be entangled have largely arisen from unsuitable names (some have even used the discordant name of impossible quantities). Had we started from the concepts offered by two-dimensional manifolds, as seen in their greatest purity in our intuitions of space, and called positive quantities direct, negative quantities inverse and imaginary quantities lateral, simplicity would have taken the place of complications, and clarity the place of darkness.

\Article{39}
What was set out in the preceding article concerns continuous complex quantities. In arithmetic, which deals only with integers, the representation of complex numbers will be a system of equally spaced points arranged on equally spaced lines so as to divide the infinite plane into infinitely many equal squares. All numbers divisible by a given complex number $a+bi=m$ will likewise form infinitely many squares with sides $=\surd(aa+bb)$, or areas $=aa+bb$. The latter squares will be inclined to the former whenever neither of the numbers $a,b$ is $=0$. Each number not divisible by the modulus $m$ will correspond to a point either inside such a square or on a side shared by two squares; the latter case, however, cannot occur unless $a,b$ have a common divisor. It is also clear that numbers congruent modulo $m$ occupy corresponding positions in their squares. Hence, if we collect all the numbers lying inside a specified square, together with all those that may lie on two of its non-opposite sides, and finally add a number divisible by $m$, we readily obtain a complete system of incongruent residues modulo $m$: every integer must be congruent to one of them, and to only one. Nor would it be difficult to show that the number of these residues equals the norm of the modulus, namely $=aa+bb$. But it seems advisable to prove this very important theorem in another, purely arithmetical way.

\Article{40}
\textsc{Theorem.}\quad\emph{For a given complex modulus $m=a+bi$, whose norm is $aa+bb=p$ and for which $a,b$ are relatively prime, every complex integer will be congruent to a residue in the sequence $0,1,2,3\dotsfour p-1$, and to no more than one.}


\SourceLeaf{121}%
\emph{Proof.} I. Let $\alpha,\beta$ be integers such that $\alpha a+\beta b=1$. Then
\[
i=\alpha b-\beta a+m(\beta+\alpha i)
\]
Thus, for a given complex integer $A+Bi$, we have
\[
A+Bi=A+(\alpha b-\beta a)B+m(\beta B+\alpha Bi)
\]
Denote by $h$ the least positive residue of $A+(\alpha b-\beta a)B$ modulo $p$, and put
\[
A+(\alpha b-\beta a)B=h+kp=h+m(ak-bki)
\]
Then
\[
A+Bi=h+m(\beta B+ak+(\alpha B-bk)i)
\]
or
\[
A+Bi\equiv h(\mathrm{mod.}\ m).\qquad\text{Q. E. P.}
\]

II. Whenever two real numbers $h,h'$ are congruent to the same complex number modulo $m$, they are also congruent to each other. We therefore put $h-h'=m(c+di)$, giving
\[
(h-h')(a-bi)=p(c+di)
\]
and hence
\[
(h-h')a=pc,\quad(h-h')b=-pd
\]
and also, since $a\alpha+b\beta=1$,
\[
h-h'=p(c\alpha-d\beta),\ \text{i. e.}\ h\equiv h'(\mathrm{mod.}\ p)
\]
Thus $h$ and $h'$, if distinct, cannot both belong to the collection of numbers $0,1,2,3\dotsfour p-1$. Q. E. S.

\newpage
\Article{41}
\textsc{Theorem.}\quad\emph{For a complex modulus $m=.\,a+bi$, whose norm is $aa+bb=p$ and for which $a,b$ are not relatively prime but have greatest common divisor $\lambda$ (which we suppose taken positively), every complex number is congruent to a residue $x+yi$ such that $x$ is one of the numbers $0,1,2,3\dotsfour\frac p\lambda-1$ and $y$ one of the numbers $0,1,2,3\dotsfour\lambda-1$, and to exactly one of all the $p$ residues of this form.}

\SourceLeaf{122}%
\emph{Proof.} I. Taking integers $\alpha,\beta$ such that $\alpha a+\beta b=\lambda$, we have $\lambda i=\alpha b-\beta a+m(\beta+\alpha i)$. Now let $A+Bi$ be a given complex number, $y$ the least positive residue of $B$ modulo $\lambda$, and $x$ the least positive residue of $A+(\alpha b-\beta a)\pdot\frac{B-y}{\lambda}$ modulo $\frac p\lambda$, and put
\[
A+(\alpha b-\beta a)\pdot\frac{B-y}{\lambda}=x+\frac p\lambda\pdot k
\]
Then
\[
\begin{aligned}
A+Bi-(x+yi)&=\frac p\lambda\pdot k+(B-y)i-(\alpha b-\beta a)\frac{B-y}{\lambda}\\
&=\frac p\lambda\pdot k+\frac{B-y}{\lambda}\pdot m(\beta+\alpha i)\\
&=\left(\frac a\lambda-\frac b\lambda\pdot i\right)km+\frac{B-y}{\lambda}(\beta+\alpha i)m
\end{aligned}
\]
which is divisible by $m$; equivalently, $A+Bi\equiv x+yi(\mathrm{mod.}\ m)$ Q. E. P.

II. Suppose that, modulo $m$, two numbers $x+yi,x'+y'i$ are congruent to the same complex number, and therefore also to each other modulo $m$. All the more, they are congruent modulo $\lambda$, so $y\equiv y'(\mathrm{mod.}\lambda)$. If both $y,y'$ are therefore assumed to belong to the numbers $0,1,2,3\dotsfour\lambda-1$, we must have $y=y'$. We then also have $x\equiv x'(\mathrm{mod.}\ m)$: $x-x'$ is divisible by $m$, and consequently the integer $\frac{x-x'}\lambda$ is divisible by $\frac a\lambda+\frac b\lambda\pdot i$, or
\[
\frac{x-x'}\lambda\equiv0\left(\mathrm{mod.}\frac a\lambda+\frac b\lambda\pdot i\right)
\]
But since $\frac a\lambda,\frac b\lambda$ are relatively prime, the second part of the theorem in the preceding article shows that $\frac{x-x'}\lambda$ is also divisible by the norm of $\frac a\lambda+\frac b\lambda\pdot i$, that is, by the number $\frac p{\lambda\lambda}$, and consequently $x-x'$ by $\frac p\lambda$. Thus, if both $x,x'$ are also assumed to belong to the collection of numbers $0,1,2,3\dotsfour\frac p\lambda-1$, we must have $x=x'$, or the residues $x+yi,x'+y'i$ are identical. Q. E. S.

It is also evident that this includes the case where the modulus is real, namely $b=0$ and consequently $\lambda=\pm a$, as well as the case where it is purely imaginary, namely $a=0$ and consequently $\lambda=\pm b$. In either case we have $\frac p\lambda=\lambda$.

\SourceLeaf{123}%
\Article{42}
Thus, assigning complex numbers congruent to one another modulo a given modulus to the same class, and incongruent numbers to different classes, we obtain altogether $p$ classes exhausting the whole range of integers, where $p$ denotes the norm of the modulus. A collection of that many numbers, one taken from each class, gives a complete system of incongruent residues, such as we specified in Articles 40, 41. In that system, the choice of residues representing their classes was based on the principle of adopting in each class a residue $x+yi$ for which $y$ has its least value and, among all those with that same least value of $y$, the one for which $x$ has its least value, excluding negative values for both $x$ and $y$. But other purposes will call for other principles. In particular, we should note the method of choosing residues that give the simplest quotients when divided by the modulus. If $\alpha+\beta i,\alpha'+\beta'i,\alpha''+\beta''i$, etc., are quotients obtained by dividing congruent numbers by the modulus, the differences among $\alpha,\alpha',\alpha''$, etc., are evidently integers, as are the differences among $\beta,\beta',\beta''$, etc. It is clear that there is always one residue for which $\alpha$ and $\beta$ lie between the bounds $0$ and $1$, the former included and the latter excluded. We simply call such a residue a least residue. If preferred, the bounds $-\frac12$ and $+\frac12$ may be adopted instead, one included and the other excluded; we shall call a residue satisfying this restriction \emph{absolutely least}.

Concerning\textperiodcentered these least residues, the following problems arise.

\Article{43}
The least residue of a given complex number $A+Bi$ modulo $a+bi$, of norm $=p$, is found as follows. If $x+yi$ is the required least residue, $(x+yi)(a-bi)$ will be the least residue of the product $(A+Bi)(a-bi)$ modulo $(a+bi)(a-bi)$, that is, modulo $p$. Put
\[
aA+bB=Fp+f,\quad aB-bA=Gp+g
\]
so that $f,g$ are the least residues of $aA+bB,aB-bA$ modulo $p$. Then
\SourceLeaf{124}%
\[
x+yi=\frac{f+gi}{a-bi}
\]
or
\[
\begin{aligned}
x&=\frac{af-bg}{p}=A-aF+bG\\
y&=\frac{ag+bf}{p}=B-aG-bF
\end{aligned}
\]
Evidently, the least residues $f,g$ must be taken either between the bounds $0$ and $p-1$ or between $-\frac12p$ and $+\frac12p$, according as the complex number's simply least or absolutely least residue is required.

\Article{44}
A complete system of least residues for a given modulus can be constructed in several ways. The first method begins by determining the bounds within which the real terms must lie, and then, for each value within those bounds, assigning bounds to the imaginary parts. The general criterion for a least residue $x+yi$ modulo $a+bi$ is that both $ax+by=\xi$ and $ay-bx=\eta$ lie between the bounds $0$ and $aa+bb$ when simply least residues are considered, or between $-\frac12(aa+bb)$ and $+\frac12(aa+bb)$ when absolutely least residues are required, one bound being excluded. Special rules would require distinguishing the cases arising from the different signs of the numbers $a,b$. Since developing that distinction presents no difficulty, we refrain from dwelling on it here; it will suffice to show the nature of the method by a single example.

\Needspace{12\baselineskip}
For the modulus $5+2i$, the simply least residues $x+yi$ must be such that both $5x+2y=\xi$ and $5y-2x=\eta$ equal one of the numbers $0,1,2,3\dotsfour28$. The equation $29x=5\xi-2\eta$ shows that positive values of $x$ cannot exceed $\frac{5\pdot28}{29}$, and negative values, disregarding their sign, cannot exceed $\frac{2\pdot28}{29}$. All admissible values of $x$ are therefore $-1,0,1,2,3,4$. For $x=-1$, $2y$ must equal one of the numbers $5,6,7\dotsfour33$, and $5y$ one of the numbers $-2,-1,0,1\dotsfour26$; hence the least value of $y$ is $+3$ and the greatest $+5$. Treating the remaining values of $x$ in the same way gives the following table of all the least residues:

\SourceLeaf{125}%
\[
\begin{array}{r|l}
x&y\\\hline
-1&3,4,5\\
0&0,1,2,3,4,5\\
+1&1,2,3\pdot4,5\pdot6\\
+2&1,2,3,4,5,6\\
+3&2,3,4\pdot5,6\\
+4&2,3,4
\end{array}
\]

\Needspace{19\baselineskip}
Similarly, for absolutely least residues, $\xi$ and $\eta$ must equal one of the numbers $-14,-13,-12\dotsfour+14$. Hence $29x$ cannot lie outside the bounds $-7\pdot14$ and $+7\pdot14$, so $x$ must equal one of the numbers $-3,-2,-1,0,1,2,3$. For $x=-3$, $2y=\xi-5x=\xi+15$ will equal one of the numbers $1,2,3\dotsfour29$, while $5y=\eta+2x=\eta-6$ will equal one of the numbers $-20,-19,-18\dotsfour+8$; this gives for $y$ the single value $+1$. Treating the remaining values of $x$ in the same way, we have the table of all the absolutely least residues:

\[
\begin{array}{r|l}
x&y\\\hline
-3&+1\\
-2&-2,-1,0,+1,+2\\
-1&-3,-2,-1,0,+1,+2\\
0&-2,-1,0,+1,+2\\
+1&-2,-1,0,+1,+2,+3\\
+2&-2,-1,0,+1,+2\\
+3&-1
\end{array}
\]

